% !TeX root = ../navier-stokes-manuscript.tex
% ============================================================
% 1.9 One Fluid, One Coherent Field, and the VPI Law
% ============================================================
\subsection{One Fluid, One Coherent Field, and the VPI Law}\label{sub:OneFluidOneCoherentFieldAndTheVPILaw}

The fluid in the problem statement is one velocity field under one simultaneous constraint.
Every point and packet is a local portion of that field.
Every derivative rung records another way that same portion responds.
The pressure at the point is fixed through the configuration of the whole field, and viscosity relates the point to the spatial curvature in its immediate neighbourhood.
These relations together are the one-fluid object.

We use \emph{VPI} as shorthand for the joined action of viscosity, pressure, and incompressibility in the velocity law.
The VPI law is the coupled system
\[
  D_tu=-\nabla p+\nu\Delta u,
  \qquad
  \nabla\cdot u=0,
  \qquad
  -\Delta p=\partial_i\partial_j(u_i u_j).
\]
The first relation gives the material response.
The second states the local volume constraint.
The third is the elliptic pressure closure forced by that constraint.
The pressure equation is derived from the first two relations, so it records their coupling and adds no separate fluid property.

This law contains two different spatial reaches at once.
Transport and viscosity are expressed through derivatives at a location.
Pressure is obtained by solving an elliptic equation across the domain.
For a sufficiently decaying field on \(\mathbb R^3\), the latter relation can be written schematically as
\[
  p(x)
  =
  \operatorname{p.v.}
  \int_{\mathbb R^3}
  K_{ij}(x-y)u_i(y)u_j(y)\,dy,
  \qquad
  K_{ij}
  =
  \partial_i\partial_j
  \frac{1}{4\pi|\cdot|},
\]
up to the usual pressure normalisation.
The periodic problem has the corresponding mean-zero Green kernel.
A change far from \(x\) can alter the pressure response at \(x\), even though the viscous term at \(x\) reads local curvature.

\divider{Carrying VPI Through the Tower}

The older and stronger meaning of the VPI law is a mathematical discipline.
Viscosity, pressure, and incompressibility must stay attached to the same transport row whenever the equation is differentiated, estimated, localised, or tested for failure.
Isolating one term in an estimate does not give that term an independent evolution.

The \(n\)-th temporal row can be collected into the residual triple
\[
  \mathscr R_n[u,p]
  :=
  \left(
    \mathfrak D_n[u,p],
    \nabla\cdot V_n,
    -\Delta P_n
    -
    \partial_i\partial_j
    \sum_{k=0}^{n}
    \binom{n}{k}
    V_{k,i}V_{n-k,j}
  \right),
\]
where \(\mathfrak D_n\) is the differentiated momentum residual and \(V_n=\partial_t^nu\), \(P_n=\partial_t^np\).
For a smooth participating Navier--Stokes field,
\[
  \mathscr R_n[u,p]=(0,0,0)
  \qquad
  \text{for every }n\geq0,
\]
and the same vanishing holds after the spatial differentiations needed to read the mixed tower.
The triple makes the participation test explicit.
A local momentum identity with the wrong divergence or pressure closure is incomplete, just as a pressure or viscous estimate detached from the velocity row is incomplete.

This point also fixes the role of a large or vanishing scalar quantity.
A large velocity gradient, a zero acceleration, or a growing Sobolev norm may warn that a particular tower rung is under stress.
The readout becomes a class-exit witness only when it identifies a failure of the joined field relation at the regularity being claimed.
The VPI discipline carries the original fluid properties through the calculation until that failure is either found or excluded.

\divider{A Graph of One Field}

Graph language can make the participation geometry visible, provided that the graph keeps the global pressure relation.
Cover the domain by overlapping patches \(\{B_a\}_{a\in\mathcal V}\).
For finite derivative depths \(N\) and \(M\), attach to each vertex \(a\) the local record
\[
  J_a^{N,M}
  :=
  \left\{
    \left.
    \partial_t^n\partial_x^\alpha u
    \right|_{B_a}
    :
    0\leq n\leq N,\ 
    |\alpha|\leq M
  \right\}.
\]
The vertices are local windows into the fluid, and \(J_a^{N,M}\) is the portion of the derivative tower visible on the \(a\)-th patch.

Define a local edge whenever two patches overlap,
\[
  \mathcal E_{\mathrm{loc}}
  :=
  \left\{
    (a,b):
    B_a\cap B_b\neq\varnothing
  \right\}.
\]
Along such an edge, the two local records must agree on their overlap,
\[
  \left.
  J_a^{N,M}
  \right|_{B_a\cap B_b}
  =
  \left.
  J_b^{N,M}
  \right|_{B_a\cap B_b}.
\]
These edges represent the compatibility needed for one velocity field and the local spatial coupling read by transport and viscosity.

Pressure requires a second relation.
Let \(\mathcal E_p\) connect patches whose velocity values are coupled by the Green kernel of the pressure equation.
In a discrete or patchwise approximation, this relation takes the form
\[
  p_a
  =
  \sum_{b\in\mathcal V}
  W_{ab}^{ij}
  (u_i u_j)_b,
\]
where \(W_{ab}^{ij}\) is the patchwise representation of the pressure kernel.
The matrix \(W\) is generally dense, so distant vertices remain related.
The exact continuum equation remains the governing object; the graph is a representation of its local and nonlocal dependencies.

At any finite depth, the resulting participation graph may be written as
\[
  \mathcal G_{N,M}
  =
  \left(
    \mathcal V,
    \mathcal E_{\mathrm{loc}},
    \mathcal E_p,
    \{J_a^{N,M}\}_{a\in\mathcal V}
  \right).
\]
It represents one coherent field only when three conditions hold together.
Overlapping records agree, every local record satisfies the required differentiated momentum and divergence rows, and every pressure value comes from the same global pressure solve.
A nearest-neighbour graph can capture overlap and local derivatives.
It cannot carry incompressible pressure by itself.

The physical picture is now complete.
One portion of the fluid changes through the velocity field that surrounds it, through viscous curvature nearby, and through the pressure constraint imposed by the whole incompressible field.
Its acceleration, jerk, and higher derivatives are successive expressions of that same response.
Restart at a regular time takes a new slice through the object without creating a new object.
Dead, Jump, and Blown identify ways in which a proposed endpoint can lose lawful response, compatible field value, or finite tower control.

VPI is the rule that keeps all of these statements about the same fluid.
Gold asks for control strong enough to keep the complete participant smooth.
Silver asks whether any nonsmooth proposal can still satisfy the complete participation law.
The two routes meet at the same class boundary because there has only ever been one fixed-viscosity, pressure-coupled, incompressible field.
